\documentclass[11pt,letterpaper]{article}

% Paquetes necesarios
\usepackage[utf8]{inputenc}
\usepackage[spanish]{babel}
\usepackage[margin=2cm]{geometry}
\usepackage{fancyhdr}
\usepackage{graphicx}
\usepackage{listings}
\usepackage{xcolor}
\usepackage{hyperref}
\usepackage{tabularx}
\usepackage{array}

% Configuración del header y footer
\pagestyle{fancy}
\fancyhf{}
\lhead{NLPY-1: Fundamentos de Python}
\rhead{Semana 6: Archivos y Persistencia}
\rfoot{Página \thepage}

% Configuración de listings para Python
\lstset{
    language=Python,
    basicstyle=\ttfamily\small,
    keywordstyle=\color{blue}\bfseries,
    stringstyle=\color{red},
    commentstyle=\color{gray}\itshape,
    numbers=left,
    numberstyle=\tiny\color{gray},
    stepnumber=1,
    numbersep=8pt,
    showstringspaces=false,
    breaklines=true,
    frame=single,
    backgroundcolor=\color{gray!10},
    tabsize=4,
    captionpos=b
}

% Configuración de hyperref
\hypersetup{
    colorlinks=true,
    linkcolor=blue,
    urlcolor=blue,
    citecolor=blue
}

\begin{document}

% Título
\begin{center}
    {\LARGE \textbf{Tarea 6: Agenda de Contactos}}\\[0.3cm]
    {\Large Semana 6: Archivos y Persistencia de Datos}\\[0.2cm]
    {\large NLPY-1: Fundamentos de Python}\\[0.2cm]
\end{center}

\section{Introducción}

En esta sexta semana del curso, exploraremos uno de los aspectos más importantes de la programación: \textbf{la persistencia de datos}. Hasta ahora, todos los programas que hemos desarrollado pierden su información cuando finalizan su ejecución. En la vida real, las aplicaciones necesitan almacenar y recuperar información de manera permanente.

El manejo de archivos es fundamental en cualquier aplicación profesional, desde sistemas de gestión empresarial hasta aplicaciones científicas que procesan grandes volúmenes de datos. Python ofrece herramientas poderosas y elegantes para trabajar con archivos de texto, archivos CSV, JSON y otros formatos.

En esta tarea desarrollarás una \textbf{Agenda de Contactos} completa que permite almacenar, buscar, modificar y eliminar información de contactos, manteniendo todos los datos guardados en archivos que persisten entre ejecuciones del programa.

\subsection{Objetivos de Aprendizaje}

Al completar esta tarea, el estudiante será capaz de:

\begin{itemize}
    \item Comprender los conceptos de persistencia de datos y operaciones de archivo
    \item Utilizar las funciones \texttt{open()}, \texttt{read()}, \texttt{write()}, y \texttt{close()} correctamente
    \item Aplicar el gestor de contexto \texttt{with} para manejo seguro de archivos
    \item Trabajar con diferentes modos de apertura de archivos (r, w, a, r+, etc.)
    \item Procesar archivos CSV utilizando el módulo \texttt{csv}
    \item Implementar operaciones CRUD (Create, Read, Update, Delete) con archivos
    \item Manejar excepciones específicas de operaciones de archivo
    \item Validar y limpiar datos antes de guardarlos
    \item Organizar código en funciones especializadas para operaciones de archivo
\end{itemize}

\section{Descripción del Proyecto}

Desarrollarás un sistema de \textbf{Agenda de Contactos} que permite gestionar información personal y profesional de contactos. El sistema debe almacenar los datos en archivos CSV y proporcionar una interfaz de menú para realizar todas las operaciones.

\subsection{Funcionalidades Requeridas}

\subsubsection{1. Gestión de Contactos (40 puntos - Funcionalidad)}

Tu programa debe implementar las siguientes funcionalidades:

\textbf{A) Agregar Nuevo Contacto (10 puntos)}
\begin{itemize}
    \item Solicitar: nombre completo, teléfono(s), email, dirección, categoría (familia/amigo/trabajo/otro)
    \item Validar formato de email (debe contener @ y dominio)
    \item Validar que el nombre no esté vacío
    \item Permitir múltiples teléfonos separados por comas
    \item Guardar en archivo CSV automáticamente
    \item Confirmar la adición exitosa
\end{itemize}

\textbf{B) Listar Todos los Contactos (5 puntos)}
\begin{itemize}
    \item Mostrar todos los contactos en formato tabular
    \item Incluir numeración para cada contacto
    \item Mostrar mensaje apropiado si no hay contactos
    \item Formato limpio y legible
\end{itemize}

\textbf{C) Buscar Contactos (8 puntos)}
\begin{itemize}
    \item Búsqueda por nombre (coincidencia parcial, case-insensitive)
    \item Búsqueda por categoría
    \item Búsqueda por teléfono
    \item Mostrar todos los resultados encontrados
    \item Indicar si no se encontraron coincidencias
\end{itemize}

\textbf{D) Actualizar Contacto (8 puntos)}
\begin{itemize}
    \item Buscar contacto a actualizar
    \item Mostrar datos actuales
    \item Permitir modificar cada campo (dejar vacío para no cambiar)
    \item Validar nuevos datos
    \item Actualizar el archivo CSV
    \item Confirmar la actualización
\end{itemize}

\textbf{E) Eliminar Contacto (5 puntos)}
\begin{itemize}
    \item Buscar y mostrar contacto a eliminar
    \item Solicitar confirmación antes de eliminar
    \item Eliminar del archivo CSV
    \item Confirmar la eliminación
\end{itemize}

\textbf{F) Exportar/Importar (4 puntos)}
\begin{itemize}
    \item Exportar contactos a archivo de texto legible
    \item Crear respaldo (backup) del archivo CSV
    \item Mostrar estadísticas (total de contactos por categoría)
\end{itemize}

\subsubsection{2. Requisitos Técnicos}

\textbf{Manejo de Archivos:}
\begin{itemize}
    \item Usar gestor de contexto \texttt{with} en TODAS las operaciones de archivo
    \item Archivo principal: \texttt{contactos.csv}
    \item Manejar correctamente el encoding UTF-8
    \item Implementar manejo de excepciones para operaciones de archivo
\end{itemize}

\textbf{Formato CSV:}
\begin{itemize}
    \item Usar módulo \texttt{csv} de Python
    \item Incluir encabezados en el archivo
    \item Campos: nombre, telefonos, email, direccion, categoria, fecha\_creacion
\end{itemize}

\textbf{Validaciones:}
\begin{itemize}
    \item Nombre: no vacío, mínimo 3 caracteres
    \item Email: formato válido (contiene @ y punto)
    \item Teléfono: solo números, guiones y espacios
    \item Categoría: debe ser una de las opciones válidas
\end{itemize}

\textbf{Organización del Código:}
\begin{itemize}
    \item Función para cargar contactos desde CSV
    \item Función para guardar contactos en CSV
    \item Función para cada operación CRUD
    \item Función de menú principal
    \item Funciones de validación separadas
\end{itemize}

\section{Ejemplo de Implementación}

A continuación se muestra un ejemplo completo del funcionamiento esperado:

\subsection{Estructura de Datos}

\begin{lstlisting}[caption={Ejemplo de archivo contactos.csv}]
nombre,telefonos,email,direccion,categoria,fecha_creacion
Juan Perez,8888-8888|7777-7777,juan@email.com,San Jose Centro,trabajo,2024-10-24
Maria Rodriguez,6666-6666,maria@email.com,Cartago,familia,2024-10-24
Carlos Mora,5555-5555,carlos@email.com,Heredia,amigo,2024-10-24
\end{lstlisting}

\subsection{Plantilla de Código}

A continuación se muestra la estructura básica que debes seguir para implementar tu solución:

\begin{lstlisting}[caption={Plantilla básica del sistema}]
import csv
import os
from datetime import datetime

# Constantes
ARCHIVO_CONTACTOS = 'contactos.csv'
CATEGORIAS_VALIDAS = ['familia', 'amigo', 'trabajo', 'otro']
ENCABEZADOS = ['nombre', 'telefonos', 'email', 'direccion', 
               'categoria', 'fecha_creacion']

def crear_archivo_si_no_existe():
    """Crea el archivo CSV con encabezados si no existe."""
    # TODO: Implementar creacion de archivo con DictWriter
    pass

def cargar_contactos():
    """Carga todos los contactos desde el archivo CSV."""
    # TODO: Usar DictReader para leer el archivo
    # TODO: Manejar FileNotFoundError
    # TODO: Retornar lista de diccionarios
    pass

def guardar_contactos(contactos):
    """Guarda todos los contactos en el archivo CSV."""
    # TODO: Usar DictWriter para escribir todos los contactos
    # TODO: Incluir encabezados
    # TODO: Manejar excepciones
    pass

def validar_email(email):
    """Valida el formato basico de un email."""
    # TODO: Verificar que contiene @ y punto en el dominio
    pass

def validar_nombre(nombre):
    """Valida que el nombre sea valido."""
    # TODO: Verificar longitud minima
    # TODO: Verificar que contiene letras
    pass

def validar_telefono(telefono):
    """Valida que el telefono solo contenga caracteres permitidos."""
    # TODO: Verificar caracteres numericos y separadores
    pass

def agregar_contacto():
    """Agrega un nuevo contacto a la agenda."""
    # TODO: Solicitar datos del contacto
    # TODO: Validar cada campo
    # TODO: Crear diccionario con los datos
    # TODO: Agregar a la lista y guardar
    pass

def listar_contactos():
    """Lista todos los contactos de la agenda."""
    # TODO: Cargar contactos
    # TODO: Mostrar en formato tabular
    # TODO: Manejar agenda vacia
    pass

def buscar_contactos():
    """Busca contactos por diferentes criterios."""
    # TODO: Mostrar menu de busqueda
    # TODO: Implementar busqueda por nombre
    # TODO: Implementar busqueda por categoria
    # TODO: Implementar busqueda por telefono
    # TODO: Mostrar resultados
    pass

def actualizar_contacto():
    """Actualiza la informacion de un contacto existente."""
    # TODO: Buscar contacto a actualizar
    # TODO: Mostrar datos actuales
    # TODO: Solicitar nuevos valores
    # TODO: Validar y actualizar
    pass

def eliminar_contacto():
    """Elimina un contacto de la agenda."""
    # TODO: Buscar contacto
    # TODO: Mostrar informacion
    # TODO: Solicitar confirmacion
    # TODO: Eliminar y guardar
    pass

def exportar_a_texto():
    """Exporta todos los contactos a un archivo de texto legible."""
    # TODO: Cargar contactos
    # TODO: Crear archivo de texto
    # TODO: Escribir con formato legible
    pass

def crear_respaldo():
    """Crea una copia de respaldo del archivo de contactos."""
    # TODO: Leer archivo original
    # TODO: Crear copia con timestamp
    pass

def mostrar_estadisticas():
    """Muestra estadisticas sobre los contactos."""
    # TODO: Contar contactos por categoria
    # TODO: Calcular porcentajes
    # TODO: Mostrar informacion resumida
    pass

def menu_principal():
    """Muestra el menu principal y maneja la navegacion."""
    crear_archivo_si_no_existe()
    
    while True:
        # TODO: Mostrar menu de opciones
        # TODO: Leer opcion del usuario
        # TODO: Llamar a la funcion correspondiente
        # TODO: Manejar opcion de salir
        pass

# Ejecutar el programa
if __name__ == "__main__":
    menu_principal()
\end{lstlisting}

\subsubsection{Ejemplo de Implementación: Función cargar\_contactos()}

Como guía, aquí hay UN EJEMPLO de cómo implementar la función de carga:

\begin{lstlisting}[caption={Ejemplo: Implementación de cargar\_contactos}]
def cargar_contactos():
    """Carga todos los contactos desde el archivo CSV."""
    contactos = []
    try:
        with open(ARCHIVO_CONTACTOS, 'r', newline='', 
                  encoding='utf-8') as archivo:
            lector = csv.DictReader(archivo)
            for fila in lector:
                contactos.append(fila)
        return contactos
    except FileNotFoundError:
        print("Archivo no encontrado. Se creara uno nuevo.")
        crear_archivo_si_no_existe()
        return []
    except Exception as e:
        print(f"Error al cargar: {e}")
        return []
\end{lstlisting}

\textbf{Nota importante}: Este es SOLO un ejemplo de referencia para una función. El resto del código debe ser implementado por el estudiante siguiendo la plantilla.

\subsection{Ejemplo de Ejecución}

\begin{lstlisting}[caption={Ejemplo de ejecución del programa}]
Bienvenido a la Agenda de Contactos
Desarrollado para NLPY-1 - Semana 6

==================================================
AGENDA DE CONTACTOS
==================================================
1. Agregar nuevo contacto
2. Listar todos los contactos
3. Buscar contactos
4. Actualizar contacto
5. Eliminar contacto
6. Exportar a archivo de texto
7. Crear respaldo
8. Mostrar estadisticas
9. Salir
==================================================

Seleccione una opcion: 1

==================================================
AGREGAR NUEVO CONTACTO
==================================================

Nombre completo: Maria Fernandez Lopez
Telefono(s) (separados por comas): 8888-8888, 2222-2222
Email: maria.fernandez@email.com
Direccion: Cartago, Dulce Nombre
Categorias disponibles: familia, amigo, trabajo, otro
Categoria: familia

¡Contacto agregado exitosamente!
Total de contactos: 4

==================================================
AGENDA DE CONTACTOS
==================================================
...
Seleccione una opcion: 3

==================================================
BUSCAR CONTACTOS
==================================================
1. Buscar por nombre
2. Buscar por categoria
3. Buscar por telefono

Seleccione tipo de busqueda: 2
Categorias: familia, amigo, trabajo, otro
Ingrese categoria: familia

2 contacto(s) encontrado(s):
--------------------------------------------------------------------------------

Nombre: Maria Rodriguez
Telefonos: 6666-6666
Email: maria@email.com
Categoria: familia

Nombre: Maria Fernandez Lopez
Telefonos: 8888-8888, 2222-2222
Email: maria.fernandez@email.com
Categoria: familia

==================================================
AGENDA DE CONTACTOS
==================================================
...
Seleccione una opcion: 8

==================================================
ESTADISTICAS DE LA AGENDA
==================================================

Total de contactos: 4

Contactos por categoria:
  Familia: 2 (50.0%)
  Amigo: 1 (25.0%)
  Trabajo: 1 (25.0%)
  Otro: 0 (0.0%)

Contacto mas antiguo: Juan Perez (2024-10-24)
\end{lstlisting}

\section{Defensa del Código (60\% de la Nota)}

La defensa oral es un componente crucial de la evaluación. Debes demostrar comprensión profunda de:

\subsection{Aspectos Técnicos a Dominar}

\textbf{1. Manejo de Archivos (15 puntos)}
\begin{itemize}
    \item ¿Qué es el gestor de contexto \texttt{with}? ¿Por qué es importante?
    \item Diferencias entre modos de apertura: 'r', 'w', 'a', 'r+', 'w+'
    \item ¿Qué pasa si no se cierra un archivo correctamente?
    \item Explicar el concepto de \textit{file handle} o descriptor de archivo
    \item ¿Cómo se maneja el encoding en archivos de texto?
\end{itemize}

\textbf{2. Módulo CSV (10 puntos)}
\begin{itemize}
    \item ¿Por qué usar el módulo \texttt{csv} en lugar de procesar texto manualmente?
    \item Diferencia entre \texttt{csv.reader} y \texttt{csv.DictReader}
    \item ¿Qué es el parámetro \texttt{newline=''} y por qué se usa?
    \item Explicar \texttt{DictWriter} y el uso de \texttt{fieldnames}
    \item ¿Cómo se manejan caracteres especiales (comas, saltos de línea) en CSV?
\end{itemize}

\textbf{3. Operaciones de Archivo (10 puntos)}
\begin{itemize}
    \item Explicar el flujo completo de agregar un contacto (desde input hasta disco)
    \item ¿Cómo funciona la actualización de archivos CSV? ¿Se puede editar "en el lugar"?
    \item Estrategias para búsqueda eficiente en archivos
    \item ¿Qué es el módulo \texttt{os} y para qué lo usamos?
    \item Diferencia entre rutas absolutas y relativas
\end{itemize}

\textbf{4. Validación y Manejo de Errores (10 puntos)}
\begin{itemize}
    \item Excepciones específicas de archivos: \texttt{FileNotFoundError}, \texttt{PermissionError}, \texttt{IOError}
    \item ¿Por qué validar datos ANTES de guardarlos?
    \item Estrategias de validación de email, teléfono, etc.
    \item ¿Qué es la sanitización de datos?
    \item Recuperación ante errores durante escritura de archivos
\end{itemize}

\textbf{5. Organización y Buenas Prácticas (15 puntos)}
\begin{itemize}
    \item Justificar la estructura del código en funciones
    \item ¿Por qué usar constantes para nombres de archivo y categorías?
    \item Explicar el patrón "cargar-modificar-guardar"
    \item Importancia de los respaldos automáticos
    \item Comentarios y documentación del código
    \item Nombres descriptivos de variables y funciones
\end{itemize}

\subsection{Preguntas de Defensa Típicas}

Prepárate para responder preguntas como:

\begin{itemize}
    \item "Muéstrame en tu código dónde usas el gestor de contexto \texttt{with}. ¿Qué pasaría si no lo usaras?"
    \item "¿Cómo maneja tu programa el caso donde el archivo de contactos no existe?"
    \item "Si el programa se interrumpe mientras se está guardando, ¿qué pasa con los datos?"
    \item "Explica línea por línea tu función de búsqueda. ¿Por qué usaste ese enfoque?"
    \item "¿Qué mejoras implementarías para manejar miles de contactos?"
    \item "Muéstrame cómo validaste el formato del email. ¿Es suficiente tu validación?"
    \item "¿Por qué elegiste CSV en lugar de otro formato como JSON o texto plano?"
    \item "Explica cómo almacenas múltiples teléfonos en un solo campo CSV"
\end{itemize}

\section{Rúbrica de Evaluación}

\subsection{Funcionalidad del Código (40 puntos)}

\textbf{Agregar Contacto (10 puntos):}
\begin{itemize}
    \item 3 puntos: Solicita y almacena todos los campos requeridos
    \item 3 puntos: Validaciones funcionan correctamente (email, nombre, teléfono)
    \item 2 puntos: Guarda correctamente en CSV con encoding UTF-8
    \item 2 puntos: Confirmación y manejo de errores
\end{itemize}

\textbf{Listar Contactos (5 puntos):}
\begin{itemize}
    \item 2 puntos: Carga y muestra todos los contactos correctamente
    \item 2 puntos: Formato tabular claro y legible
    \item 1 punto: Maneja caso de agenda vacía
\end{itemize}

\textbf{Buscar Contactos (8 puntos):}
\begin{itemize}
    \item 3 puntos: Búsqueda por nombre funciona (parcial, case-insensitive)
    \item 2 puntos: Búsqueda por categoría funciona
    \item 2 puntos: Búsqueda por teléfono funciona
    \item 1 punto: Mensajes apropiados cuando no hay resultados
\end{itemize}

\textbf{Actualizar Contacto (8 puntos):}
\begin{itemize}
    \item 2 puntos: Encuentra y muestra contacto correctamente
    \item 3 puntos: Permite modificar campos con validación
    \item 2 puntos: Guarda cambios correctamente en CSV
    \item 1 punto: Maneja múltiples contactos con nombre similar
\end{itemize}

\textbf{Eliminar Contacto (5 puntos):}
\begin{itemize}
    \item 2 puntos: Encuentra y muestra contacto a eliminar
    \item 2 puntos: Solicita confirmación y elimina correctamente
    \item 1 punto: Actualiza archivo CSV
\end{itemize}

\textbf{Exportar/Estadísticas (4 puntos):}
\begin{itemize}
    \item 2 puntos: Exporta a texto legible correctamente
    \item 1 punto: Crea respaldos
    \item 1 punto: Muestra estadísticas por categoría
\end{itemize}

\subsection{Defensa del Código (60 puntos)}

\textbf{Comprensión de Archivos (15 puntos):}
\begin{itemize}
    \item 5 puntos: Explica gestor de contexto y modos de apertura
    \item 5 puntos: Demuestra comprensión de lectura/escritura
    \item 5 puntos: Entiende encoding y manejo de caracteres especiales
\end{itemize}

\textbf{Módulo CSV (10 puntos):}
\begin{itemize}
    \item 5 puntos: Explica uso de \texttt{DictReader} y \texttt{DictWriter}
    \item 3 puntos: Comprende manejo de delimitadores y formato
    \item 2 puntos: Explica ventajas sobre procesamiento manual
\end{itemize}

\textbf{Operaciones CRUD (10 puntos):}
\begin{itemize}
    \item 3 puntos: Explica flujo de creación de contactos
    \item 3 puntos: Comprende lectura y búsqueda
    \item 2 puntos: Explica actualización correctamente
    \item 2 puntos: Comprende eliminación segura
\end{itemize}

\textbf{Validación y Errores (10 puntos):}
\begin{itemize}
    \item 4 puntos: Identifica y maneja excepciones de archivo
    \item 3 puntos: Explica validaciones implementadas
    \item 3 puntos: Demuestra comprensión de sanitización de datos
\end{itemize}

\textbf{Organización y Calidad (15 puntos):}
\begin{itemize}
    \item 5 puntos: Código bien organizado en funciones
    \item 3 puntos: Uso apropiado de constantes
    \item 3 puntos: Nombres descriptivos y comentarios útiles
    \item 2 puntos: Manejo del módulo \texttt{datetime}
    \item 2 puntos: Implementa buenas prácticas de respaldo
\end{itemize}

\subsection{Penalizaciones}

\begin{itemize}
    \item \textbf{-5 puntos}: No usar gestor de contexto \texttt{with} para archivos
    \item \textbf{-5 puntos}: No manejar encoding UTF-8 correctamente
    \item \textbf{-3 puntos}: No validar datos antes de guardar
    \item \textbf{-3 puntos}: Código sin organización en funciones
    \item \textbf{-3 puntos}: No usar el módulo \texttt{csv} (procesar manualmente)
    \item \textbf{-2 puntos}: No manejar excepciones de archivo
    \item \textbf{-2 puntos}: Variables o funciones con nombres no descriptivos
    \item \textbf{-2 puntos}: No incluir comentarios explicativos
    \item \textbf{-5 puntos}: Plagio o código copiado sin comprensión
\end{itemize}

\section{Desafíos Opcionales (Puntos Extra)}

\subsection{Nivel Intermedio (+10 puntos cada uno)}

\textbf{1. Importar desde CSV externo}
\begin{itemize}
    \item Permitir importar contactos desde otro archivo CSV
    \item Validar formato y detectar duplicados
    \item Opción de fusionar o reemplazar contactos existentes
\end{itemize}

\textbf{2. Exportar a formato JSON}
\begin{itemize}
    \item Implementar exportación a formato JSON
    \item Usar módulo \texttt{json}
    \item Mantener estructura de datos adecuada
\end{itemize}

\textbf{3. Búsqueda avanzada}
\begin{itemize}
    \item Búsqueda con múltiples criterios simultáneos
    \item Operadores lógicos (AND, OR)
    \item Ordenar resultados por diferentes campos
\end{itemize}

\subsection{Nivel Avanzado (+15 puntos cada uno)}

\textbf{4. Sistema de registro de cambios (log)}
\begin{itemize}
    \item Crear archivo de log que registre todas las operaciones
    \item Incluir timestamp, usuario (nombre del equipo), acción realizada
    \item Usar módulo \texttt{logging}
\end{itemize}

\textbf{5. Manejo de fotografías}
\begin{itemize}
    \item Permitir asociar una foto a cada contacto
    \item Almacenar ruta de la imagen en el CSV
    \item Validar que el archivo de imagen exista
    \item Crear carpeta de imágenes organizada
\end{itemize}

\textbf{6. Cifrado básico de datos}
\begin{itemize}
    \item Implementar cifrado simple de emails o teléfonos sensibles
    \item Usar módulo \texttt{base64} o similar
    \item Agregar opción de cifrar/descifrar
\end{itemize}

\subsection{Nivel Experto (+20 puntos)}

\textbf{7. Sincronización con múltiples archivos}
\begin{itemize}
    \item Mantener versiones de respaldo automáticas
    \item Detectar y resolver conflictos de datos
    \item Implementar sistema de versionado
\end{itemize}

\textbf{8. Interfaz gráfica básica}
\begin{itemize}
    \item Crear interfaz gráfica simple con \texttt{tkinter}
    \item Mantener todas las funcionalidades del sistema
    \item Mejorar experiencia de usuario
\end{itemize}

\section{Consejos y Recomendaciones}

\subsection{Para el Desarrollo}

\begin{enumerate}
    \item \textbf{Comienza simple}: Implementa primero agregar y listar, luego añade funcionalidades
    \item \textbf{Prueba frecuentemente}: Verifica cada función antes de continuar
    \item \textbf{Usa datos de prueba}: Crea contactos variados para probar todos los casos
    \item \textbf{Maneja errores}: Anticipa qué puede salir mal y manéjalo
    \item \textbf{Versiona tu código}: Guarda versiones de respaldo mientras desarrollas
    \item \textbf{Documenta mientras programas}: No dejes la documentación para el final
\end{enumerate}

\subsection{Para la Defensa}

\begin{enumerate}
    \item \textbf{Conoce tu código}: Entiende cada línea que escribiste
    \item \textbf{Practica explicaciones}: Ensaya cómo explicarías conceptos clave
    \item \textbf{Prepara el ambiente}: Ten tu código listo para ejecutar y mostrar
    \item \textbf{Anticipa preguntas}: Piensa qué preguntarías tú sobre tu código
    \item \textbf{Sé honesto}: Si no sabes algo, reconócelo y muestra disposición a aprender
    \item \textbf{Muestra ejemplos}: Ten casos de prueba preparados para demostrar funcionalidad
\end{enumerate}

\subsection{Recursos de Apoyo}

\begin{itemize}
    \item Documentación oficial de Python: \url{https://docs.python.org/3/}
    \item Tutorial de archivos: \url{https://docs.python.org/3/tutorial/inputoutput.html#reading-and-writing-files}
    \item Módulo CSV: \url{https://docs.python.org/3/library/csv.html}
    \item Módulo OS: \url{https://docs.python.org/3/library/os.html}
    \item Módulo datetime: \url{https://docs.python.org/3/library/datetime.html}
\end{itemize}

\section{Criterios de Entrega}

\subsection{Archivos a Entregar}

\begin{enumerate}
    \item \textbf{agenda\_contactos.py}: Código principal del programa
    \item \textbf{README.txt}: Instrucciones de uso y descripción
    \item \textbf{contactos.csv}: Archivo de ejemplo con al menos 5 contactos de prueba
    \item \textbf{requirements.txt}: Dependencias (si usas librerías adicionales)
\end{enumerate}

\subsection{Formato del README.txt}

Tu archivo README debe incluir:

\begin{itemize}
    \item Nombre del estudiante
    \item Descripción breve del proyecto
    \item Instrucciones de ejecución
    \item Funcionalidades implementadas
    \item Desafíos opcionales completados (si aplica)
    \item Problemas conocidos o limitaciones
    \item Recursos utilizados para el desarrollo
\end{itemize}

\subsection{Evaluación de la Defensa}

La defensa será individual y consistirá en:

\begin{enumerate}
    \item \textbf{Demostración (5 minutos)}: Ejecutar el programa y mostrar funcionalidades
    \item \textbf{Explicación de código (10 minutos)}: Explicar secciones específicas del código
    \item \textbf{Preguntas conceptuales (5 minutos)}: Responder preguntas sobre conceptos teóricos
    \item \textbf{Resolución de problemas (5 minutos)}: Modificar el código para agregar una pequeña funcionalidad
\end{enumerate}

\section{Notas Finales}

\begin{itemize}
    \item \textbf{Fecha de entrega}: Según calendario del curso
    \item \textbf{Peso de la tarea}: 10\% de la nota final del curso
    \item \textbf{Nota mínima}: 60/100 para aprobar esta tarea
    \item \textbf{Consultas}: Usar los canales oficiales del curso
    \item \textbf{Integridad académica}: El código debe ser 100\% propio. Plagios serán penalizados con 0.
\end{itemize}

\vspace{1cm}

\begin{center}
\textit{"Los datos son el nuevo petróleo, pero a diferencia del petróleo,}\\
\textit{los datos no se agotan cuando se usan... si sabes cómo almacenarlos correctamente."}\\
\vspace{0.3cm}
--- Adaptación libre sobre la importancia de la persistencia de datos
\end{center}

\vspace{0.5cm}

\begin{center}
\textbf{¡Éxito en tu desarrollo!}\\
\textit{Recuerda: Un buen sistema de archivos es la base de cualquier aplicación seria.}
\end{center}

\end{document}
